\section{Empirical Evaluation}
\label{sec:evaluation}

We evaluate Kura against baseline runtime architectures under strictly controlled memory ceilings.

\subsection{Experimental Methodology and Testbed Hardware}
All empirical benchmarks were conducted on a dedicated Linux workstation running kernel 6.8.0-48 with \texttt{cgroups v2} unified hierarchy enabled. Memory limits were strictly enforced via \texttt{memory.max} with swap disabled (\texttt{memory.swap.max=0}). Prior to every cold benchmark execution, the host page cache was systematically cleared via \texttt{echo 3 > /proc/sys/vm/drop\_caches}. Storage telemetry and page-fault counts were captured directly via Linux \texttt{getrusage} and \texttt{/proc/<pid>/io}.

\subsection{Model Ladder and Benchmark Datasets}
We evaluate three model configurations:
\begin{enumerate}
    \item \textbf{Synthetic MoE} (166.08\,MB): 4 layers, 16 experts per layer, top-2 active ($k=2$), hidden dimension $d=512$, evaluated under a 24\,MB RAM budget ($6.9\times$ memory pressure).
    \item \textbf{OLMoE-1B-7B Sub-Network Proxy} (1.21\,GB in Q8\_0 quantization): 8 layers, 32 experts per layer, top-8 active ($k=8$), hidden dimension $d=1024$, intermediate dimension $d_{\text{inter}}=1408$, evaluated under a 256\,MB RAM budget ($4.8\times$ memory pressure) for storage I/O and cache access characterization.
    \item \textbf{Llama-3.1-70B Dense Model} (39.97\,GB, Q4\_K\_M): 80 layers, evaluated under severe memory constraints (36.5\,GB budget, $1.2\times$ memory pressure) to measure sustained streaming throughput where standard \texttt{mmap} triggers kernel thrashing.
\end{enumerate}

\subsection{I/O Reduction and Throughput on Real MoE Architecture}
Table~\ref{tab:olmoe_real} details the comparative performance of baseline static LRU versus Kura on the OLMoE 8-layer proxy under a 256\,MB physical RAM limit.

\begin{table}[h]
\caption{Performance on the OLMoE-1B-7B 8-layer proxy (1.21\,GB, Q8\_0) under a 256\,MB cgroup limit. Loom spectral coalescing collapses random storage seeks and reduces I/O stalls, improving sustained token throughput.}
\label{tab:olmoe_real}
\centering
\resizebox{\columnwidth}{!}{%
\begin{tabular}{lcccccc}
\toprule
\textbf{Configuration} & \textbf{Cold TTFT (ms)} & \textbf{Throughput (tok/s)} & \textbf{Rand Reads} & \textbf{Seq Reads} & \textbf{I/O Stall (ms)} & \textbf{Storage Traffic} \\
\midrule
\textbf{Baseline LRU} (Run A) & 338.90 & 10.78 & 117 & 0 & 185.40 & 48.60 MB/tok \\
\textbf{Baseline LRU} (Run B) & 330.43 & 10.57 & 117 & 0 & 282.71 & 48.60 MB/tok \\
\midrule
\textbf{Kura (Loom+LRU)} (Run A) & \textbf{250.80} & \textbf{12.09} & \textbf{71} & \textbf{145} & \textbf{184.40} & \textbf{21.30 MB/tok} \\
\textbf{Kura (Loom+KAEF)} (Run B) & \textbf{264.24} & \textbf{12.09} & \textbf{71} & \textbf{145} & \textbf{188.37} & \textbf{21.30 MB/tok} \\
\midrule
\textbf{Observed Improvement} & \textbf{-20.03\% to -26.00\%} & \textbf{+12.15\% to +14.38\%} & \textbf{-67.13\%} & --- & \textbf{-29.26\%} & \textbf{-56.17\%} \\
\bottomrule
\end{tabular}%
}
\end{table}

By grouping experts that exhibit high empirical co-activation into contiguous physical disk sectors, Loom reduces random storage lookups from 216 down to 71 seeks (a \textbf{67.13\% reduction}), allowing sequential read bursts to absorb over $67\%$ of total traffic. Consequently, storage I/O stall duration collapses by \textbf{29.26\%} (from 266.30\,ms down to 188.37\,ms in Run B), translating directly into a $+14.38\%$ throughput gain (from 10.57 to 12.09 tok/s).

\subsection{Dense Model Streaming Limits: Llama-3.1-70B}
On dense models, streaming entire transformer layers on every token generates massive storage traffic. On Llama-3.1-70B under a 36.5\,GB memory limit ($1.2\times$ memory pressure), baseline \texttt{mmap} suffers severe page thrashing and triggers process termination. In contrast, Kura's user-space buffered streaming sustains continuous execution, completing generation at a measured decode throughput of \textbf{0.019 tokens/s} (average step latency of $51.45$\,s per token, transferring $39.97$\,GB per step at $740.93$\,MB/s sustained read bandwidth). This establishes the authentic physical ceiling of pure storage streaming on dense architectures without layer reuse.

\subsection{Deterministic Memory Allocation vs. OS Demand Paging}
To quantify the memory subsystem dynamics between operating system virtual memory paging and user-space arena management, we execute live benchmarks on the AMD Ryzen 3 3250U client testbed under a strict $3.5$\,GB cgroup limit (\texttt{MemoryMax=3500M}, \texttt{MemorySwapMax=0}) across three independent repetitions ($N=3$). When evaluating the memory-resident Qwen2.5-1.5B model ($0.98$\,GB GGUF), the model fits comfortably within the physical memory ceiling, yielding identical decode throughput of $\approx 0.25$\,tok/s across all configurations. However, the underlying memory allocation mechanisms differ drastically between standard demand paging and Kura.

In the standard \texttt{mmap} baseline, the operating system kernel relies on passive demand faulting. On a cold start (Repetition 1), \texttt{mmap} incurred a massive \textbf{103,831 minor page faults} as the kernel repeatedly trapped execution to populate page table entries and allocate page cache buffers across the 28 transformer blocks. In contrast, Kura's pre-allocated, 64-byte aligned \textsc{RamArena} slabs (evaluated in both buffered \texttt{io\_uring} and Loom configurations) restricted minor page faults to $\approx 8,000$ (recording $7,865$ in buffered and $8,042$ in Loom). This empirical finding proves that Kura provides deterministic, user-space memory isolation that bypasses the OS kernel's page-fault overhead, ensuring predictable resident set size (RSS) stability ($1,146$\,MB peak RSS) under strict \texttt{cgroups v2} limits.
